%%%%%%%%%%%%%%%%%%%%%%%%%%%%%%%%%%%%%%%%%%%%%%%%%%%%%%%
% NOTAS TEORIA DE TRANSPORTE OPTIMO
% FCEN-UBA, segundo cuatrimestre de 2026. Fuente unica de las notas
% y de las guias de ejercicios (herramientas/generar_practicas.py).
%%%%%%%%%%%%%%%%%%%%%%%%%%%%%%%%%%%%%%%%%%%%%%%%%%%%%%%

\documentclass[11pt,a4paper]{amsart}
\usepackage{xr-hyper}
\newcounter{chapter}
\usepackage{fancyhdr}
\usepackage{geometry} \geometry{top=2cm,bottom=2.5cm,left=2.5cm,right=2.5cm}

\usepackage{amssymb}
\usepackage{hyperref}
\usepackage{listings}
\usepackage{mathtools}
\usepackage{enumitem}

% Fuzz -------------------------------------------------------------------
% Line spacing -----------------------------------------------------------

%%%%%%%%%fin pegada%%%%%%%%%%%%%%%%%%5

\usepackage[T1]{fontenc}
\usepackage[spanish]{babel}
\usepackage{graphicx}
\usepackage{tikz}
\usetikzlibrary{babel}



\parskip 3pt

%%
%% OPERADORES
%%
\renewcommand{\Re}{\operatorname{Re}}
\newcommand{\tr}{\operatorname{tr}}
\newcommand{\id}{\operatorname{id}}
\def\div{\operatorname {\mathrm{div}}}
\def\Dom{\operatorname {\mathrm{Dom}}}
\def\argmin{\operatorname {\mathrm{arg\, min}}}
\def\argmax{\operatorname {\mathrm{arg\, max}}}
\def\gen{\operatorname {\mathrm{gen}}}
\def\Im{\operatorname {\mathrm{Im}}}
\def\Nu{\operatorname {\mathrm{Nu}}}
\def\sop{\operatorname {\mathrm{sop}}}
\def\diam{\operatorname {\mathrm{diam}}}
\def\dist{\operatorname {\mathrm{dist}}}
\def\var{\operatorname {\mathrm{Var}}}
\def\codim{\operatorname {\mathrm{codim}}}
\def\pint{\operatorname {--\!\!\!\!\!\int\!\!\!\!\!--}}

%%
%% EPSILON
%%
\newcommand{\ve}{\varepsilon}

%%
%% LETRAS MATHBB
%%
\newcommand{\R}{{\mathbb R}}
\newcommand{\Z}{{\mathbb Z}}
\newcommand{\C}{{\mathbb C}}
\newcommand{\Q}{{\mathbb Q}}
\newcommand{\N}{{\mathbb N}}
\newcommand{\E}{{\mathbb E}}
\newcommand{\V}{{\mathbb V}}
\renewcommand{\P}{{\mathbb P}}
\newcommand{\II}{{\mathbb I}}
\newcommand{\OO}{{\mathbb O}}
\newcommand{\Sb}{{\mathbb S}}
\newcommand{\Tb}{{\mathbb T}}

%%
%% LETRAS CALIFGRAFICAS
%%
\newcommand{\F}{{\mathcal F}}
\newcommand{\Hc}{{\mathcal H}}
\newcommand{\Tc}{{\mathcal T}}
\newcommand{\Ec}{{\mathcal E}}
\newcommand{\Vc}{{\mathcal V}}
\newcommand{\A}{{\mathcal A}}
\newcommand{\J}{{\mathcal J}}
\newcommand{\G}{{\mathcal G}}
\newcommand{\U}{{\mathcal U}}
\newcommand{\B}{{\mathcal B}}
\newcommand{\Sw}{{\mathcal S}}
\newcommand{\D}{{\mathcal D}}
\newcommand{\M}{{\mathcal M}}
\newcommand{\LL}{{\mathcal L}}

%%
%% LETRAS MATHBF
%%
\newcommand{\ab}{{\mathbf a}}
\newcommand{\bb}{{\mathbf b}}
\newcommand{\cc}{{\mathbf c}}
\newcommand{\ee}{{\mathbf e}}
\newcommand{\ff}{{\mathbf f}}
\newcommand{\gb}{{\mathbf g}}
\newcommand{\hh}{{\mathbf h}}
\newcommand{\n}{{\mathbf n}}
\newcommand{\mm}{{\mathbf m}}
\newcommand{\pp}{{\mathbf p}}
\newcommand{\qq}{{\mathbf q}}
\newcommand{\sbf}{{\mathbf s}}
\newcommand{\tb}{{\mathbf t}}
\newcommand{\uu}{{\mathbf u}}
\newcommand{\vv}{{\mathbf v}}
\newcommand{\ww}{{\mathbf w}}
\newcommand{\xx}{{\mathbf x}}
\newcommand{\yy}{{\mathbf y}}
\newcommand{\zz}{{\mathbf z}}
\newcommand{\I}{{\mathbf 1}}
\newcommand{\Pb}{{\mathbf P}}
\newcommand{\XX}{{\mathbf X}}
\newcommand{\YY}{{\mathbf Y}}

%%
%% LETRAS MATHFRAK
%%

\newcommand{\Gk}{{\mathfrak G}}


%%
%% TEOREMAS
%%
\newtheorem{teo}{Teorema}[section]
\newtheorem{lema}[teo]{Lema}
\newtheorem{prop}[teo]{Proposición}
\newtheorem{cor}[teo]{Corolario}

%%
%% REMARK
%%
\theoremstyle{remark}
\newtheorem{obs}[teo]{Observación}

%%
%% DEFINICION
%%
\theoremstyle{definition}
\newtheorem{defi}[teo]{Definición}
\newtheorem{ejem}[teo]{Ejemplo}
\newtheorem{ejer}[teo]{Ejercicio}
\newtheorem{nota}[teo]{Notación}
\newtheorem{ejercicio}{Ejercicio}[chapter]


%%
%% NUMERACION
%%
\numberwithin{subsection}{section}
\numberwithin{equation}{section}

%%
%% TITULO
%%


\vfuzz7pt \hfuzz7pt
\externaldocument[N-]{../notas/Notas-TO}
\newcommand{\encabezado}[1]{%
\noindent \includegraphics[scale=2.1]{logo-dm-wide.png} \hskip 9.5cm
\includegraphics[scale=0.08]{logo-exactas-uba.jpg}
\smallskip
\usefont{T1}{phv}{m}{n}
\begin{center}{\Large\sc Teoría de Transporte Óptimo}\end{center}
\bigskip
\noindent{\sc #1}\\[-0.4cm]
\hrule
\normalfont}
\newcommand{\cuadernolink}[2]{\noindent\textbf{Cuaderno:} \emph{#1} (\href{https://colab.research.google.com/github/jfbonder/transporte-optimo/blob/main/cuadernos/#2}{abrir en Colab}).\par\medskip}

\begin{document}
\renewcommand{\thechapter}{13}
\encabezado{Práctica 13: Transporte entrópico y el algoritmo de Sinkhorn}
\bigskip
\pagestyle{fancy}
\fancyhead[L]{{\footnotesize{\usefont{T1}{phv}{m}{n}{\sc Teoría de Transporte Óptimo}}}}
\fancyhead[R]{\footnotesize \usefont{T1}{phv}{m}{n}{\sc Práctica 13}}
\thispagestyle{plain}
\noindent{\footnotesize Los ejercicios son los de la sección de ejercicios del Capítulo~13 de las notas del curso, con la misma numeración; las referencias a teoremas, proposiciones y secciones remiten a las notas (\url{https://github.com/jfbonder/transporte-optimo}).}
\medskip

\cuadernolink{Transporte entrópico y Sinkhorn}{06-transporte-entropico-sinkhorn.ipynb}


\begin{ejercicio}[Propiedades de la entropía relativa]\label{ejer.KL.propiedades}
Sean $a\in\R^n$, $b\in\R^m$ con entradas positivas y $\pi\in\Pi(a,b)$.
\begin{enumerate}[label=(\alph*)]
\item Probar que $\mathrm{KL}(\pi\,|\,a\otimes b)=\sum_{ij}\pi_{ij}\log\pi_{ij}-\sum_ia_i\log a_i-\sum_jb_j\log b_j$;
es decir, sobre $\Pi(a,b)$ la entropía relativa difiere de la entropía
negativa $\sum\pi_{ij}\log\pi_{ij}$ en una constante. Deducir que el
regularizador de Peyré--Cuturi de la Observación~\ref{N-obs.notacion.peyre}
produce el mismo minimizador $\pi^\ve$.
\item (Desigualdad de Pinsker.) Probar que
$\|\pi-a\otimes b\|_1^2\le2\,\mathrm{KL}(\pi\,|\,a\otimes b)$.
\emph{Sugerencia:} reducirlo al caso de dos puntos (agrupar según el signo
de $\pi_{ij}-a_ib_j$ y usar la desigualdad de log-suma
$\sum_ks_k\log\frac{s_k}{r_k}\ge\bigl(\sum_ks_k\bigr)\log\frac{\sum_ks_k}{\sum_kr_k}$,
que es la convexidad conjunta de $(s,r)\mapsto s\log\frac sr$) y estudiar la
función $t\mapsto t\log\frac tp+(1-t)\log\frac{1-t}{1-p}-2(t-p)^2$.
\item Deducir de (b) y de la Proposición~\ref{N-prop.eps.infinito} la
velocidad $\|\pi^\ve-a\otimes b\|_1\le2\sqrt{\|C\|_\infty/\ve}$.
\item Probar que la cota $\mathrm{KL}(\pi\,|\,a\otimes b)\le\log\frac1{\min_ia_i}$
usada en el Teorema~\ref{N-teo.eps.cero} es óptima: para $a=b$ uniformes en
$n$ puntos y $\pi$ una matriz de permutación dividida por $n$ hay igualdad.
\end{enumerate}
\end{ejercicio}

\begin{ejercicio}[El valor entrópico como función de $\ve$]\label{ejer.OTeps.concavo}
\leavevmode
\begin{enumerate}[label=(\alph*)]
\item Probar que $\ve\mapsto\mathrm{OT}_\ve(a,b)$ es no decreciente y cóncava
en $(0,\infty)$ (es un ínfimo de funciones afines de $\ve$), y que
\[
\frac{d}{d\ve}\mathrm{OT}_\ve(a,b)=\mathrm{KL}(\pi^\ve\,|\,a\otimes b)
\]
donde sea derivable.
\item Deducir que $\ve\mapsto\mathrm{KL}(\pi^\ve\,|\,a\otimes b)$ es no
creciente y que $\ve\mapsto\langle C,\pi^\ve\rangle$ es no decreciente:
al aumentar la regularización el plan se vuelve más «desordenado» y más
caro.
\item Probar que $\lim_{\ve\to\infty}\mathrm{OT}_\ve(a,b)=\langle C,a\otimes b\rangle$
y que $\lim_{\ve\to0}\mathrm{OT}_\ve(a,b)=\mathrm{OT}(a,b)$, y que ambas
convergencias son monótonas.
\end{enumerate}
\end{ejercicio}

\begin{ejercicio}[El soft-min]\label{ejer.softmin}
Con la notación de la Definición~\ref{N-def.softmin}:
\begin{enumerate}[label=(\alph*)]
\item Probar que $(\varphi+s)^{c,\ve}=(\varphi)^{c,\ve}-s$ para $s\in\R$, y
que $\varphi\le\tilde\varphi$ implica $(\varphi)^{c,\ve}\ge(\tilde\varphi)^{c,\ve}$.
\item Probar que $\varphi\mapsto(\varphi)^{c,\ve}_j$ es cóncava y de clase
$C^\infty$, y calcular su gradiente: es (menos) el vector de probabilidad
$\bigl(a_ie^{(\varphi_i-c_{ij})/\ve}/\sum_ka_ke^{(\varphi_k-c_{kj})/\ve}\bigr)_i$,
la $j$-ésima columna normalizada del plan. Comparar con la $c$-transformada
$\varphi^c_j=\min_i\{c_{ij}-\varphi_i\}$, que no es diferenciable.
\item Probar que $\varphi\mapsto(\varphi)^{c,\ve}$ es $1$-Lipschitz para la
norma $\|\cdot\|_\infty$, y que para $\ve\to\infty$,
$(\varphi)^{c,\ve}_j\to\sum_ia_i(c_{ij}-\varphi_i)$: el soft-min tiende al
promedio.
\end{enumerate}
\end{ejercicio}

\begin{ejercicio}[Unicidad de los potenciales entrópicos]\label{ejer.potenciales.eps.unicidad}
Sean $(\varphi,\psi)$ y $(\varphi',\psi')$ dos pares tales que
$a_ib_je^{(\varphi_i+\psi_j-c_{ij})/\ve}$ y $a_ib_je^{(\varphi'_i+\psi'_j-c_{ij})/\ve}$
son ambos planes de $\Pi(a,b)$.
\begin{enumerate}[label=(\alph*)]
\item Probar que ambos planes coinciden (son el único minimizador
$\pi^\ve$), y deducir que $\varphi_i+\psi_j=\varphi'_i+\psi'_j$ para todo
$(i,j)$.
\item Concluir que $(\varphi',\psi')=(\varphi+s,\psi-s)$ para algún $s\in\R$.
Contrastar con el caso $\ve=0$ (Ejercicio~\ref{N-ejer.discreto.unicidad}), en
el que los potenciales pueden no ser únicos: ¿qué propiedad del plan
entrópico lo explica?
\item Probar que si $a=b$ y $C$ es simétrica, entonces $\varphi^\ve=\psi^\ve$
con la normalización adecuada, y que en ese caso una iteración de Sinkhorn
puede escribirse con un solo potencial.
\end{enumerate}
\end{ejercicio}

\begin{ejercicio}[Los potenciales entrópicos convergen a potenciales de Kantorovich]\label{ejer.potenciales.eps.limite}
Completar los detalles de la Observación~\ref{N-obs.potenciales.eps}.
\begin{enumerate}[label=(\alph*)]
\item Normalizando $\varphi^\ve_1=0$, probar que $(\varphi^\ve,\psi^\ve)$
está acotado uniformemente para $\ve\in(0,1]$. \emph{Sugerencia:}
$\psi^\ve=(\varphi^\ve)^{c,\ve}$ y el Lema~\ref{N-lema.softmin} dan
$\|\psi^\ve-(\varphi^\ve)^c\|_\infty\le\ve\log\frac1{\min a_i}$; usar que
$\varphi^\ve=(\psi^\ve)^{\bar c,\ve}$ tiene oscilación a lo sumo
$\max_{i,i',j}|c_{ij}-c_{i'j}|$ (Paso~2 de la demostración del
Teorema~\ref{N-teo.sinkhorn.convergencia}), de modo que, con
$\varphi^\ve_1=0$, $\varphi^\ve$ y luego $(\varphi^\ve)^c$ están acotadas
independientemente de $\ve$.
\item Probar que todo punto límite $(\varphi^0,\psi^0)$ satisface
$\psi^0=(\varphi^0)^c$, $\varphi^0=(\psi^0)^{\bar c}$ y
$\varphi^0_i+\psi^0_j=c_{ij}$ sobre el soporte de $\pi^0$, y concluir que es
un par óptimo del problema dual de Kantorovich.
\item ¿Converge toda la familia $(\varphi^\ve,\psi^\ve)$, o solo
subsucesiones? Discutirlo en el ejemplo del
Ejercicio~\ref{N-ejer.discreto.3x3}, donde los potenciales de Kantorovich
forman un segmento: ¿a qué punto del segmento converge el límite entrópico?
(La respuesta es delicada; ver \cite{Cominetti1994}.)
\end{enumerate}
\end{ejercicio}

\begin{ejercicio}[Sinkhorn en un paso]\label{ejer.sinkhorn.un.paso}
\leavevmode
\begin{enumerate}[label=(\alph*)]
\item Probar que si $C$ es de la forma $c_{ij}=f_i+g_j$ entonces
$\pi^\ve=a\otimes b$ para todo $\ve$, y el algoritmo de Sinkhorn converge en
una iteración.
\item Recíprocamente, probar que si $\pi^\ve=a\otimes b$ para algún
$\ve>0$, entonces $C$ es de la forma (a). Deducir que $\pi^\ve$ tiene rango
uno si y solo si $C$ es de la forma (a).
\item Sea $n=m=2$. Escribir explícitamente $\pi^\ve$ en función de $a$, $b$,
$C$ y $\ve$ (el problema se reduce a una ecuación escalar) y verificar a
mano los límites $\ve\to0$ y $\ve\to\infty$.
\end{enumerate}
\end{ejercicio}

\begin{ejercicio}[Baricentros entrópicos]\label{ejer.bar.entropico}
\leavevmode
\begin{enumerate}[label=(\alph*)]
\item Probar directamente que $a\mapsto\widetilde{\mathrm{OT}}_\ve(a,b^k)$ es
estrictamente convexa en $\Sigma_n$, sin usar la entropía $H(a)$:
si $\pi$ y $\pi'$ son los planes óptimos para $a\ne a'$, el plan
$\frac12(\pi+\pi')$ es admisible para $\frac12(a+a')$, y la convexidad
estricta de $\pi\mapsto\sum\pi_{ij}(\log\pi_{ij}-1)$ da una desigualdad
estricta (¿por qué $\pi\ne\pi'$?). Deducir otra vez la unicidad del
baricentro entrópico.
\item Interpretar el paso 3 del algoritmo de la
Sección~\ref{N-sec.baricentros.entropicos} como una \emph{proyección de
Bregman}: dados los planes $\pi^1,\dots,\pi^N$ (con primeras marginales
$m^k=\pi^k\mathbf 1$ distintas), probar que el problema
\[
\min\Bigl\{\sum_k\lambda_k\,\widetilde{\mathrm{KL}}(\tilde\pi^k\,|\,\pi^k):\ \tilde\pi^k\mathbf 1=a\ \text{para todo }k,\ a\in\R^n_{\ge0}\Bigr\},
\]
donde $\widetilde{\mathrm{KL}}(p\,|\,q)=\sum p\log\frac pq-p+q$,
tiene solución $\tilde\pi^k=\operatorname{diag}(a/m^k)\,\pi^k$ con
$a=\prod_k(m^k)^{\lambda_k}$ (media geométrica ponderada, sin normalizar),
que es la actualización del algoritmo. Probar que si en cambio se usa la
$\mathrm{KL}$ usual y se impone $a\in\Sigma_n$, se obtiene la versión
normalizada $a\propto\prod_k(m^k)^{\lambda_k}$. ¿Por qué no aparece la media
aritmética?
\item En el ejemplo de no unicidad de la Observación~\ref{N-obs.bar.no.unicidad}
($\mu_1$ con átomos en $(\pm1,0)$, $\mu_2$ con átomos en $(0,\pm1)$, sobre
una grilla que contenga los cuatro candidatos $(\pm\frac12,\pm\frac12)$),
probar que el baricentro entrópico es único para todo $\ve>0$ y simétrico
respecto de los dos ejes, y describir su límite cuando $\ve\to0$.
\end{enumerate}
\end{ejercicio}

\begin{ejercicio}[Derivadas del costo de transporte]\label{ejer.OTeps.derivadas}
\leavevmode
\begin{enumerate}[label=(\alph*)]
\item \emph{Sin regularizar.} Escribamos $\mathrm{OT}_C(a,b)$ para indicar la
dependencia en la matriz de costos. Probar que, para $a,b$ fijos,
$C\mapsto\mathrm{OT}_C(a,b)$ es cóncava y que, si $\pi^0$ es un plan óptimo para $C$,
entonces $\mathrm{OT}_{C'}(a,b)\le\mathrm{OT}_C(a,b)+\langle C'-C,\pi^0\rangle$
para toda $C'$. Probar que si el plan óptimo para $C$ es único, la función es
diferenciable en $C$ con derivadas parciales $\pi^0_{ij}$. Para $n=m=2$,
$a=b=(\frac12,\frac12)$, mostrar que
$\mathrm{OT}_C(a,b)=\frac12\min\{c_{11}+c_{22},\,c_{12}+c_{21}\}$ y deducir que
no es diferenciable en $C=0$.
\item \emph{Respecto de los pesos.} Probar, sin usar la
Proposición~\ref{N-prop.baricentro.entropico} y escribiendo
$\mathrm{OT}_\ve(\cdot,b)$ como supremo de funciones afines, que, para $b$
fijo, $a\mapsto\mathrm{OT}_\ve(a,b)$ es convexa en el conjunto de vectores de
probabilidad con entradas positivas. Probar que si $h\in\R^n$ satisface
$\sum_ih_i=0$, entonces
\[
\frac{d}{dt}\Big|_{t=0}\mathrm{OT}_\ve(a+th,b)=\sum_{i=1}^nh_i\varphi^\ve_i,
\]
es decir, el potencial entrópico $\varphi^\ve$ es el gradiente de
$\mathrm{OT}_\ve$ respecto de $a$, salvo una constante (que no se ve en las
direcciones con $\sum h_i=0$). \emph{Sugerencia:} el funcional dual
$\mathcal D_\ve$ de \eqref{N-eq:dual.entropico} es afín en $a$ para
$(\varphi,\psi)$ fijos; evaluarlo en los potenciales de $a$ y en los de
$a+th$, como en el Paso 1 de la demostración de la
Proposición~\ref{N-prop.OTeps.derivada}, y usar que $\sum_j\pi^\ve_{ij}=a_i$.
\end{enumerate}
\end{ejercicio}

\begin{ejercicio}[El sesgo entrópico]\label{ejer.sinkhorn.sesgo}
Con la notación de la Observación~\ref{N-obs.divergencia.sinkhorn}, sea $c$
un costo con $c(x,x)=0$ y $c(x,y)>0$ para $x\ne y$, y sea
$\alpha=\sum_{i=1}^na_i\delta_{x_i}$ con $n\ge2$ puntos distintos y $a_i>0$.
\begin{enumerate}[label=(\alph*)]
\item Probar que $0<\mathrm{OT}_\ve(\alpha,\alpha)\le\ve H(a)$, donde
$H(a)=-\sum_ia_i\log a_i$. \emph{Sugerencia:} para la cota superior,
comparar con el plan diagonal; para la inferior, usar que el plan entrópico
tiene todas sus entradas positivas.
\item Sean $\alpha=\frac12(\delta_{-1}+\delta_1)$ y
$\beta_t=\frac12(\delta_{-t}+\delta_t)$ en $\R$, con $c(x,y)=|x-y|^2$.
Usando la Proposición~\ref{N-prop.OTeps.derivada} (para la regla de la
cadena basta la primera parte: las entradas de $C$ dependen de $t$) y la
simetría del problema, probar que
$\frac{d}{dt}\mathrm{OT}_\ve(\alpha,\beta_t)>0$ en $t=1$: el mínimo de
$\beta\mapsto\mathrm{OT}_\ve(\alpha,\beta)$ no se alcanza en $\beta=\alpha$,
y achicar el soporte de $\beta$ disminuye el costo.
\item Probar que la divergencia de Sinkhorn $S_\ve$ no cambia si en su
definición se reemplaza $\mathrm{OT}_\ve$ por el costo
$\widetilde{\mathrm{OT}}_\ve$ con el regularizador de Peyré--Cuturi
(Sección~\ref{N-sec.baricentros.entropicos}).
\item Probar que $S_\ve(\alpha,\beta)\to\mathrm{OT}(\alpha,\beta)$ cuando
$\ve\to0$, y que cuando $\ve\to\infty$
\[
S_\ve(\alpha,\beta)\to\sum_{i,j}a_ib_j\,c(x_i,y_j)
-\frac12\sum_{i,k}a_ia_k\,c(x_i,x_k)-\frac12\sum_{j,l}b_jb_l\,c(y_j,y_l).
\]
\end{enumerate}
\end{ejercicio}

\subsection*{Ejercicios computacionales}

Los ejercicios siguientes se resuelven con el cuaderno \emph{Transporte
entrópico y Sinkhorn} (véase el Apéndice~\ref{N-ap.cuadernos}); los enunciados
figuran también al final del cuaderno. Las referencias a secciones son a las
del cuaderno.

\begin{ejercicio}[Convergencia de valores]\label{ejc.eps.valores}
Graficar $\ve\mapsto\mathrm{OT}_\ve(a,b)-\mathrm{OT}(a,b)$ en escala
logarítmica en ambos ejes para el ejemplo de la Sección~1 del cuaderno y
estimar el orden de convergencia. ¿Es $O(\ve)$, como garantiza el
Teorema~\ref{N-teo.eps.cero}, o mejor? Comparar la constante observada con
$\log(1/\min_ia_i)$.
\end{ejercicio}

\begin{ejercicio}[Sinkhorn y la mejora de potenciales]\label{ejc.eps.mejora}
Implementar la mejora de potenciales $\varphi\mapsto\varphi^{c\bar c}$ de la
Proposición~\ref{N-prop:mejora-potenciales}, sin regularización, y compararla
con la iteración de Sinkhorn para $\ve$ pequeño partiendo del mismo $\varphi$.
¿Converge la versión no regularizada? ¿Y la regularizada? Comparar con el
Ejercicio~\ref{N-ejer.mejora.discreta}.
\end{ejercicio}

\begin{ejercicio}[Costo de cada iteración]\label{ejc.eps.costo}
Medir el tiempo de una iteración de Sinkhorn para $n=m=500,1000,2000,4000$ y
verificar el crecimiento $O(nm)$. Comparar con el tiempo del método exacto
(\texttt{ot.emd}) para los mismos tamaños.
\end{ejercicio}

\begin{ejercicio}[Un baricentro con $\ve\to0$]\label{ejc.eps.baricentro}
En el ejemplo de la Sección~5.1 del cuaderno, calcular el baricentro
entrópico de la Proposición~\ref{N-prop.baricentro.entropico} para
\[
\ve\in\{0{,}5,\;0{,}2,\;0{,}1,\;0{,}05,\;0{,}02,\;0{,}01,\;0{,}005\}
\]
y graficar la distancia $W_2$ entre el baricentro entrópico y el exacto en
función de $\ve$, en escala logarítmica en ambos ejes. ¿Cuántas iteraciones
hacen falta en cada caso?
\end{ejercicio}

\begin{ejercicio}[Unicidad del baricentro entrópico]\label{ejc.eps.unicidad}
Verificar numéricamente que el baricentro entrópico del ejemplo de no
unicidad de la Observación~\ref{N-obs.bar.no.unicidad} ($\mu_1$ con átomos en $(\pm1,0)$,
$\mu_2$ con átomos en $(0,\pm1)$, sobre una grilla que contenga los cuatro
candidatos) es \emph{único} y simétrico: la entropía rompe la degeneración
eligiendo la combinación más «repartida». Comparar con el
Ejercicio~\ref{ejer.bar.entropico}.
\end{ejercicio}

\bibliographystyle{plain}
\bibliography{../notas/biblio}

\end{document}
